\documentclass[aspectratio=169,11pt]{beamer}

% common_setting.tex
% Shared preamble for the five PSPDE2026 workshop beamer decks (00-04).
% \input this right after \documentclass in each deck.
% Requires xelatex (fontspec/unicode-math need XeTeX or LuaTeX, not pdflatex).

% Fonts
\usepackage{fontspec}
\usepackage{unicode-math}
\setmainfont{Liberation Sans}
\setsansfont{Liberation Sans}
\setmonofont{Liberation Mono}[Scale=0.85]
\setmathfont{Latin Modern Math}

% Theme and colors
\usetheme{moloch}
\molochset{
  block=fill,
  progressbar=frametitle,
  sectionpage=progressbar,
}
\setbeamertemplate{page number in head/foot}[totalframenumber]

% Graphics
\usepackage{graphicx}
\graphicspath{{./figures/}}

% Tables
\usepackage{booktabs}
\usepackage{tabularx}

% Unnumbered captions without the "Figure:" label prefix
\usepackage{caption}
\captionsetup{labelformat=empty}

% Code listings
\usepackage{listings}
\usepackage{xcolor}
\definecolor{codebg}{RGB}{245,245,245}
% moloch's own accent colors, aliased under their original Metropolis names
\colorlet{mAlertFg}{mLightBrown}
\colorlet{mExampleFg}{mLightGreen}
\lstdefinestyle{python}{
  language=Python,
  basicstyle=\ttfamily\scriptsize,
  keywordstyle=\color{mAlertFg}\bfseries,
  commentstyle=\color{mExampleFg}\itshape,
  stringstyle=\color{mExampleFg!70!black},
  backgroundcolor=\color{codebg},
  rulecolor=\color{black!15},
  frame=single,
  breaklines=true,
  showstringspaces=false,
  columns=fullflexible,
  numbers=none
}
\lstdefinestyle{shell}{
  basicstyle=\ttfamily\scriptsize,
  backgroundcolor=\color{codebg},
  rulecolor=\color{black!15},
  frame=single,
  breaklines=true,
  columns=fullflexible
}


% Title information
\title[BYOD]{Bring Your Own Data}
\subtitle{Block 4 --- Modern Workflow in Geology}
\author[O.\ Lexa]{Ondrej Lexa}
\institute[Charles University]{
  Institute of Petrology and Structural Geology (IPSG)\\
  Charles University, Prague, Czechia
}
\date{8 October 2026}

\begin{document}

\begin{frame}[plain]
  \titlepage
\end{frame}

\begin{frame}{This Block (15:15--16:30)}
  \begin{block}{Objective}
    Load and analyse data from your own thesis --- with time for
    consultations and debugging.
  \end{block}
  \vspace{0.3cm}
  \begin{columns}[T]
    \begin{column}{0.48\textwidth}
      \textbf{\texttt{04a\_byod\_structural.ipynb}}
      \small
      \begin{itemize}
        \item Field measurements (CSV/Excel, long or wide table)
        \item Stereonet per structure + summary with fold axis
        \item Statistics table; optional fault stress inversion
      \end{itemize}
    \end{column}
    \begin{column}{0.48\textwidth}
      \textbf{\texttt{04b\_byod\_epma.ipynb}}
      \small
      \begin{itemize}
        \item Microprobe export (Excel/CSV)
        \item APFU, end-members, QC per mineral
        \item Ternary/scatter/profile charts, one results workbook
      \end{itemize}
    \end{column}
  \end{columns}
\end{frame}

\section{Using the Templates}

\begin{frame}{Three Steps}
  \begin{enumerate}
    \item \textbf{Copy} the template (keep the original as reference) and your data
          file into \texttt{notebooks/data/}
    \item Edit \textbf{only the \texttt{CONFIG} cell}: file name, sheet, column names,
          which structures/minerals to process
    \item \emph{Run $\rightarrow$ Run All Cells} and look in \texttt{notebooks/output/}
  \end{enumerate}
  \vspace{0.4cm}
  \begin{alertblock}{Golden rule}
    Never edit the raw data file to ``make it work''. Fix things in code
    (\texttt{RENAME}, \texttt{SKIPROWS}, filters) so the path from raw data to
    figure stays reproducible.
  \end{alertblock}
\end{frame}

\begin{frame}[fragile]{Structural Template --- \texttt{CONFIG}}
  \begin{lstlisting}[style=python]
PROJECT = "field"                              # used in output file names
FILE = Path("data/field_measurements.csv")     # .csv, .xlsx or .xls
LAYOUT = "long"          # "long": structure column + azi, inc | "wide": S1_azi, S1_inc, ...
STRUCTURE_COL = "structure"
AZI_COL, INC_COL = "azi", "inc"
PLANE_NOTATION = "dd"    # "dd" dip direction/dip | "rhr" right-hand-rule strike/dip
PLANES = ["S0", "S1"]
LINES = ["L1"]
FOLDED = "S0"            # surface whose pi-girdle gives the fold axis (or None)
FAULTS = Path("data/mele.csv")  # fazi, finc, lazi, linc, sense (or None)
FAULT_SENSE_LETTERS = False     # True if sense is coded n/r/s/d instead of 1/-1
  \end{lstlisting}
  \small Output: \texttt{<PROJECT>\_<structure>.pdf}, \texttt{\_summary.pdf},
  \texttt{\_faults\_*.pdf}, \texttt{\_apsg.xlsx} (statistics, data, stress inversion)
\end{frame}

\begin{frame}[fragile]{EPMA Template --- \texttt{CONFIG}}
  \begin{lstlisting}[style=python]
SAMPLE = "S12"
FILE = Path("data/epma_export.xlsx")
SHEET = "Spots"
SKIPROWS = 0                        # header lines above the column names
RENAME = {}                         # e.g. {"SiO2(Mass%)": "SiO2"}
MINERAL_COLUMN = None               # a column that already holds the mineral ...
MINERAL_FROM = ("Comment", r"_([a-z]+)_")   # ... or a regex applied to a column
MINERALS = {"grt": (Grt, (98.5, 101.5)),    # label -> (mineral, total window)
            "pl":  (Fsp, (98.5, 101.5)),
            "bt":  (Bt,  (94.0, 97.5))}
MIN_QUALITY = 0.95                  # stoichiometry_quality threshold
PROFILE = ("Grt-profile", "Distance_um")    # optional traverse sheet, or None
  \end{lstlisting}
  \small Output: QC summary, \texttt{<SAMPLE>\_*.pdf} charts, \texttt{<SAMPLE>\_petropandas.xlsx}
\end{frame}

\section{Common Pitfalls}

\begin{frame}{Structural Data}
  \small
  \begin{tabularx}{\textwidth}{l>{\raggedright\arraybackslash}X}
    \toprule
    \textbf{Symptom} & \textbf{Likely cause / fix} \\
    \midrule
    Planes plot 90\textdegree{} off & strike/dip given as dip direction/dip --- set
      \texttt{PLANE\_NOTATION = "rhr"} (or quadrant strings: \texttt{fol("N30E,40NW")}) \\
    Poles where planes should be & plotted with \texttt{s.line} instead of
      \texttt{s.pole} / \texttt{s.great\_circle} \\
    Fold axis nonsense & mixed limbs of several fold generations, or a cluster instead
      of a girdle --- check \texttt{P}/\texttt{G} in the statistics table \\
    Stress inversion unstable & sense coding: \texttt{1} normal, \texttt{-1} reverse
      (letters \texttt{n/r/s/d}: set \texttt{FAULT\_SENSE\_LETTERS = True}); mixed
      fault populations --- filter by misfit \\
    Rows silently missing & empty cells --- accessors skip rows with NaN \\
    \bottomrule
  \end{tabularx}
\end{frame}

\begin{frame}{Microprobe Data}
  \small
  \begin{tabularx}{\textwidth}{l>{\raggedright\arraybackslash}X}
    \toprule
    \textbf{Symptom} & \textbf{Likely cause / fix} \\
    \midrule
    Oxide missing from \texttt{df.oxides()} & odd header (\texttt{SiO2 wt\%}, \texttt{SiO2(Mass\%)})
      --- add to \texttt{RENAME}; \texttt{FeO*}, \texttt{FeOT}, \texttt{H2O+} work as-is \\
    Header is garbage & instrument writes lines above the table --- set \texttt{SKIPROWS} \\
    Numbers read as text & decimal comma or ``\texttt{<LOD}'' strings ---
      \texttt{pd.read\_csv(..., decimal=",")}, \texttt{pd.to\_numeric(errors="coerce")} \\
    Fe\textsuperscript{3+} looks wrong & both \texttt{FeO} and \texttt{Fe2O3} present, or
      wrong mineral object (\texttt{Grt} vs.\ \texttt{GrtFe3}) \\
    Everything rejected & total window too tight for hydrous minerals, or wrong
      mineral label mapping in \texttt{MINERALS} \\
    \bottomrule
  \end{tabularx}
\end{frame}

\section{Wrap-up}

\begin{frame}{Where to Go Next}
  \begin{itemize}
    \item \textbf{Documentation \& tutorial notebooks}
          \begin{itemize}
            \item \texttt{github.com/ondrolexa/apsg}
            \item \texttt{github.com/ondrolexa/petropandas}
          \end{itemize}
    \item \textbf{Keep your project reproducible}: one \texttt{uv} project per thesis,
          raw data read-only, notebooks under version control (git)
    \item \textbf{Grow the templates}: they are ordinary notebooks --- add cells, cite
          the code in your methods section
    \item \textbf{Bugs and wishes}: GitHub issues on either repository
  \end{itemize}
  \vspace{0.5cm}
  \centering \Large Thank you --- now, your data!
\end{frame}

\end{document}
